\begin{Problem}

The BCS wavefunction in momentum space can be written as
\begin{equation}
|\Psi_{\mathrm{BCS}}\rangle
= \prod_{\mathbf{k}} \exp\left( \phi_{\mathbf{k}} c^\dagger_{\mathbf{k}\uparrow} c^\dagger_{-\mathbf{k}\downarrow} \right)|0\rangle
= \prod_{\mathbf{k}} \left(1 + \phi_{\mathbf{k}} c^\dagger_{\mathbf{k}\uparrow} c^\dagger_{-\mathbf{k}\downarrow}\right)|0\rangle .
\end{equation}

\begin{itemize}
\item Show that the BCS ground state is the vacuum for the Bogoliubov quasiparticles, i.e. that the destruction operators $\alpha_{\mathbf{k},\sigma}$ annihilate the BCS ground state. Here
\begin{equation}
\alpha_{\pm\mathbf{k},\sigma} = u_{\mathbf{k}} c_{\pm\mathbf{k}\sigma} + v_{\mathbf{k}} c^\dagger_{\mp\mathbf{k}\overline{\sigma}},
\end{equation}
and you may assume that $u_{\mathbf{k}}, v_{\mathbf{k}}$ are real for convenience.

\item If the Bogoliubov quasiparticle is $\alpha^\dagger_{\mathbf{k}\uparrow} = u_{\mathbf{k}} c^\dagger_{\mathbf{k}\uparrow} + v_{\mathbf{k}} c_{-\mathbf{k}\downarrow}$, then starting from the equation of motion of the Bogoliubov quasiparticle
\begin{equation}
	[H,\alpha^\dagger_{\mathbf{k}\uparrow}] = \pdv{\alpha_{\mathbf{k}\uparrow}^\dagger}{\tau}=E_{\mathbf{k}} \alpha^\dagger_{\mathbf{k}\uparrow},
\end{equation}
show that $(u_{\mathbf{k}}, v_{\mathbf{k}})^T$ is an eigenvector of
\begin{equation}
	h_{\mathbf{k}} = \begin{pmatrix} u_{\mathbf{k}} \\ v_{\mathbf{k}} \end{pmatrix} = \begin{pmatrix} \epsilon_{\mathbf{k}} & \Delta \\ \Delta & -\epsilon_{\mathbf{k}} \end{pmatrix} \begin{pmatrix} u_{\mathbf{k}} \\ v_{\mathbf{k}} \end{pmatrix} = E_{\mathbf{k}} \begin{pmatrix} u_{\mathbf{k}} \\ v_{\mathbf{k}} \end{pmatrix}  
\end{equation}
By solving the eigenvalue problem and assuming the gap is real, give  an explicit expression for $u_{\mathbf{k}}$ and $v_{\mathbf{k}}$ in terms of $\Delta$ and $\epsilon_{\mathbf{k}}$. 
\end{itemize}
\end{Problem}
\begin{proof}
	We apply the annihilation operators directly. For concreteness, we will consider $\alpha_{\va k,\uparrow}$. Since we can commute $\alpha$ through all bilinears $1+\phi_{\va q}c_{\va q}^\dagger c_{\va q}^\dagger$ with $\va q \neq \va k$, we will ignore them all and simply look at the product
	\[( u_{\mathbf{k}} c_{\mathbf{k},\uparrow} + v_{\mathbf{k}} c^\dagger_{-\mathbf{k},\downarrow})(1+\phi_{\va k}c_{\va k,\uparrow}^\dagger c_{-\va k,\downarrow}^\dagger)(1+\phi_{-\va k}c_{-\va k,\uparrow}^\dagger c_{\va k,\downarrow}^\dagger)\]
	applied to the ground state. This product would yield $2^3=8$ terms. However, we can immediately see that some fo them drop out. The product $c_{\va k,\uparrow}$ with the two 1s would yield an annihilation operator applied to the vacuum; it vanishes. We may not multiply $c_{-\va k,\downarrow}^\dagger$ with itself; squares of fermion operators vanish. Thus, two more terms drop out. If we multiply $c_{\va k,\uparrow}$ with something not containing $c_{\va k,\uparrow}^\dagger$, we could simply anticommute it to the end, where it would annihilate the ground state again. Thus, there are only 4 terms left, and we find
	\begin{align*}
		&~( u_{\va{k}} c_{\va{k},\uparrow} + v_{\va{k}} c^\dagger_{-\va{k},\downarrow})(1+\phi_{\va k}c_{\va k,\uparrow}^\dagger c_{-\va k,\downarrow}^\dagger)(1+\phi_{-\va k}c_{-\va k,\uparrow}^\dagger c_{\va k,\downarrow}^\dagger)\\
		=&~u_{\va k}c_{\va k,\uparrow} \phi_{\va k}c_{\va k,\uparrow}^\dagger c_{-\va k,\downarrow}^\dagger + u_{\va k}c_{\va k,\uparrow} \phi_{\va k}c_{\va k,\uparrow}^\dagger c_{-\va k,\downarrow}^\dagger \phi_{-\va k}c^{\dagger}_{-\va k,\uparrow}c_{\va k,\downarrow}^\dagger \\
		&+ v_{\va k}c_{-\va k,\downarrow}^\dagger + v_{\va k}\phi_{-\va k}c_{-\va k,\downarrow}^\dagger \phi_{-\va k}c_{-\va k,\uparrow}^\dagger c_{\va k,\downarrow}^\dagger\\
		=& -u_{\va k}\phi_{\va k}c_{-\va k,\downarrow}^\dagger - u_{\va k}\phi_{\va k}\phi_{-\va k}c_{-\va k,\downarrow}^\dagger c_{-\va k,\uparrow}^\dagger c_{\va k,\downarrow}^\dagger\\
		&+ v_{\va k}c_{-\va k,\downarrow}^\dagger + v_{\va k}\phi_{-\va k}c_{-\va k,\downarrow}^\dagger \phi_{-\va k}c_{-\va k,\uparrow}^\dagger c_{\va k,\downarrow}^\dagger\\
		=&\phi_{-\va k}c_{-\va k\downarrow}^\dagger c_{-\va ,\uparrow}^\dagger c_{\va k,\downarrow}^\dagger (-u_{\va k}\phi_{\va k} +v_{\va k})+c_{-\va k,\downarrow}(-u_{\va k}\phi_{\va k} +v_{\va k})
	\end{align*}
	where we have used the fermionic anticommutation relations to swap $c_{\va k,\uparrow}$ with $c_{\va k,\uparrow}^\dagger$ before pulling it through the rest of the operators to annihilate the ground state. Thus, we see that these operators annihilate the ground state if we have
	\[v_{\va k} = u_{\va k}\phi_{\va k}\]
	In principle, we could use this and the requirement that the transformation stay canonical to determine the parameters $u_{\va k}$ and $v_{\va k}$. But we continue...
	
	Next, we compute the commutator $[H, \alpha_{\va k\uparrow}^\dagger]$ directly. We have previously done a mean field expansion on the BCS Hamiltonian
	\[H = \sum_{\va k,\sigma}\epsilon_{\va k}  c_{\va k,\sigma}^\dagger c_{\va k,\sigma} -\sum_{\va k, \va k'}\frac{g}{V} c_{\va k\uparrow}^\dagger c_{-\va k\downarrow}^\dagger c_{\va k',\uparrow} c_{-\va k', \downarrow} \]
	using the mean field $\Delta = -\frac{g}{V}\langle c_{\va k,\uparrow}c_{-\va k,\downarrow}\rangle$. After decomposing, we find that
	\begin{align*}
		H &= \sum_{\va k,\sigma}\epsilon_{\va k}  c_{\va k,\sigma}^\dagger c_{\va k,\sigma} -\sum_{\va k, \va k'}\frac{g}{V} c_{\va k\uparrow}^\dagger c_{-\va k\downarrow}^\dagger c_{\va k',\uparrow} c_{-\va k', \downarrow} \\
		&= \sum_{\va k,\sigma}\epsilon_{\va k}  c_{\va k,\sigma}^\dagger c_{\va k,\sigma} -\sum_{\va k, \va k'}\frac{g}{V} \left(\overline{\Delta} - (\overline{\Delta}-c_{\va k\uparrow}^\dagger c_{-\va k\downarrow}^\dagger)\right)\left(\Delta - (\Delta - c_{\va k',\uparrow} c_{-\va k', \downarrow})\right)\\
		&= \sum_{\va k,\sigma}\epsilon_{\va k}  c_{\va k,\sigma}^\dagger c_{\va k,\sigma} -\sum_{\va k, \va k'}()
	\end{align*}
\end{proof}
\begin{Problem}

One useful way to regard superconductors is via the semiconductor analogy, in which the quasiparticles are treated like the positive and negative energy excitations of a semiconductor. To do so we divide the Brillouin zone up into equal halves and redefine a set of positive and negative energy quasiparticle operators according to
$\alpha^\dagger_{\mathbf{k},\sigma+} = \alpha^\dagger_{\mathbf{k}\sigma}$ and $\alpha^\dagger_{\mathbf{k}\sigma-} = \operatorname{sgn}(\sigma)\,a_{-\mathbf{k}-\sigma}
\text{ for } \mathbf{k} \in \tfrac12 \mathrm{BZ}.$

\begin{itemize}
\item By this procedure rewrite the BCS Hamiltonian in terms of these new operators, and show that the excitation spectrum can be interpreted in terms of an empty band of positive energy excitation and a filled band of negative energy excitations.

\item Show that the BCS ground state wave functions can be regarded as a filled sea of negative energy quasiparticle states and an empty sea of positive energy quasiparticle states.
\end{itemize}
\end{Problem}


\begin{Problem}

\begin{itemize}
\item Assuming that the Debye frequency is a small fraction of the bandwidth, show that the difference between the superconducting and normal-state free energies can be written as
\begin{equation}
\mathcal{F}_S - \mathcal{F}_N =
-2T N(0) \int_{-\omega_D}^{\omega_D}
\ln\!\left[
\frac{\cosh\!\left(\sqrt{\varepsilon^2+|\Delta|^2}/2T\right)}
{\cosh(\varepsilon/2T)}
\right] d\varepsilon
+ V \frac{|\Delta|^2}{g_0}.
\end{equation}

Notice that the previous expression is invariant under changes in the phase of the gap parameter. What consequences does it imply?

\item By differentiating the expression of the Free energy with respect to $\Delta$, confirm the zero temperature gap equation
\begin{equation}
\frac{V}{g_0 N(0)} =
\int_0^{\omega_D} \frac{d\varepsilon}{\sqrt{\varepsilon^2+\Delta_0^2}},
\end{equation}
where $\Delta_0 = \Delta(T = 0)$. Use this result to eliminate $g_0$, to show that the free energy can be written as
\begin{equation}
\mathcal{F}_S - \mathcal{F}_N =
N(0)\Delta_0^2\,
\Phi\!\left[\frac{\Delta}{\Delta_0},\frac{T}{\Delta_0}\right],
\end{equation}
where the dimensionless function
\begin{equation}
\Phi[\delta,t] =
\int_0^\infty dx
\left\{
-4t \ln
\frac{\cosh\!\left(\sqrt{x^2+\delta^2}/2t\right)}
{\cosh(x/2t)}
+ \frac{\delta^2}{\sqrt{x^2+1}}
\right\}.
\end{equation}

What implies the change in the integration boundary to infinity?
\end{itemize}
\end{Problem}


\begin{Problem}

Consider two superconductors which are connected via a small bridge, also superconducting, of length $L$. Each of those superconductors can be thought as a lead with all of its properties well defined and each of the wave function in this bridge
\begin{equation}
\Psi(x \to 0) = \Psi_1 = |\Psi_0|e^{i\phi_1}
\end{equation}
\begin{equation}
\Psi(x \to L) = \Psi_2 = |\Psi_0|e^{i\phi_2}
\end{equation}
for the edges. In this area the wave function can be thought as a superposition of those two
\begin{equation}
\Psi(x) = |\Psi_0|e^{i\phi_1} g(x) + |\Psi_0|e^{i\phi_2}(1-g(x)),
\end{equation}
with the only requirement that $g(x \to 0) \to 1$ and $g(x \to L) \to 0$. We aim to get the behaviour of $g(x)$, in order to do so use the Ginzburg Landau equation
\begin{equation}
-\frac{\hbar^2}{2m}\frac{\partial^2\Psi}{\partial x^2}
+ \alpha \Psi + \beta|\Psi|^2\Psi = 0,
\end{equation}
to have a complete description.

\textit{Hint:} It might be useful to think of $\Psi(x)=|\Psi_0|f(x)$ and rewrite an equation for $f(x)$.
Use $|\Psi_0|^2=-\alpha/\beta$ and define also $\xi^2=\hbar^2/(2m|\alpha|)$. Use the hypothesis of short link $L \ll \xi$.
\end{Problem}

\begin{Problem}
With the previous result, compute the Josephson current. Show that it depends exclusively on $\phi =\phi_2-\phi_1$ and not on each phase individually. What is the expression for the free energy?	
\end{Problem}

\begin{Problem}
	How is the current changed when one applies a voltage? Consider the Schr\"{o}dinger equation and assume $\Psi(x,t) = |\Psi(x)|e^{i\phi(t)} $. Use that to find $\phi(t)$ assuming $\phi(t=0)=\phi_0$. You may associate the energy of the external superconductors with their chemical potential, and the difference of chemical potential with the applied voltage.	
\end{Problem}
